% Part: history
% Chapter: biographies
% Section: gerhard-gentzen

\documentclass[../../../include/open-logic-section]{subfiles}

\begin{document}

\olfileid{his}{bio}{gen}

\olsection{غيرهارد غنتسن}

\olphoto{gentzen-gerhard}{Gerhard Gentzen}

اشتهر غيرهارد غنتسن بتأسيس نظرية البرهان البنيوية، ولا سيما بوضعه نظامي
!!{derivation} للاستنباط الطبيعي وحساب التعاقب. وُلد في غرايفسفالت
بألمانيا يوم \OLClassicalDigits{24} نوفمبر \OLClassicalDigits{1909}. وتعلّم في منزله ثلاث سنين قبل دخول المدرسة
الإعدادية، فكان يوم التحق بها دون أكثر أقرانه في التحصيل؛ ثم بان نبوغه
وقوة ملكته الرياضية. وكانت له عناية بأشياء شتى؛ فمن ذلك أنه كان ينظم
الشعر لأمه، ويكتب مسرحيات لمسرح المدرسة.

استهل دراسته الجامعية في غرايفسفالت، ثم تنقل بين غوتنغن وميونيخ
وبرلين. ونال الدكتوراه من جامعة غوتنغن سنة \OLClassicalDigits{1933} بإشراف هرمان فايل.
(وكان بول برنايس قد أشرف على أكثر العمل، غير أن النازيين أقصوه عن
الجامعة.) ثم صار غنتسن سنة \OLClassicalDigits{1934} مساعدًا لديفيد هيلبرت، وفي السنة نفسها
عرض نظامي !!{derivation} لحساب التعاقب والاستنباط الطبيعي في بحثيه
\emph{Untersuchungen \"{u}ber das logische Schlie\ss en I--II
  [دراسات في الاستنباط المنطقي I--II]}. وأقام سنة \OLClassicalDigits{1936} برهان اتساق
مسلمات بيانو.

وعلاقته بالنازية ملتبسة. ففي الوقت الذي أُكره أستاذه برنايس على مغادرة
ألمانيا، دخل غنتسن الفرع الجامعي من كتيبة العاصفة، وهي منظمة نازية شبه
عسكرية، وكان ككثير من الألمان عضوًا في الحزب النازي. وخدم في الحرب ضابط
اتصالات بوحدة الاستخبارات الجوية، ثم سُرّح سنة \OLClassicalDigits{1942} لانهيار عصبي. ولا
يُعلم أكان مخلصًا للحزب، أم إنما انضم إليه صونًا لمستقبله الجامعي.

وعُرض عليه سنة \OLClassicalDigits{1943} منصب في معهد الرياضيات بالجامعة الألمانية في براغ
فقبله. ثم ثار أهل براغ سنة \OLClassicalDigits{1945} على الاحتلال الألماني، ولما دخلت القوات
السوفيتية المدينة اعتقلت أساتذة الجامعة كلهم. ونُقل غنتسن، لعضويته في
منظمات نازية، إلى معسكر للسخرة؛ فمات بسوء التغذية في زنزانته يوم \OLClassicalDigits{4}
أغسطس \OLClassicalDigits{1945}، وله \OLClassicalDigits{35} سنة.

\begin{reading}
لترجمة غنتسن المستوفاة انظر \citet{Menzler-Trott2007}. وفي
\citet{Segal2014} قراءة نافعة عن الرياضيين في ظل الحكم النازي، وفيه
نبذة من حياته. وأبحاثه في الاستنباط المنطقي متاحة بأصلها الألماني
\citep{Gentzen1935a,Gentzen1935b}. وجمع \citet{Gentzen1969} ترجماتها
الإنجليزية في مجلد واحد، وألحق به ترجمة موجزة لصاحبها.
\end{reading}

\end{document}
